\section{Implementation}
\subsection{Main Functions}

\begin{itemize}
\item Simulate SP and CBS scheduling. The simulator runs for a user-specified duration (default: 4000000.0 $\mu s$) separately for SP and CBS. The maximum observed end-to-end latency for each stream within the simulation horizon is recorded as the simulated delay result.

\item Compute the analytical WCDs for both schedulers based on the theoretical formulas. For CBS, the analytical model is applied only to AVB traffic (Class A and Class B). Best Effort (BE) traffic is reported as \texttt{NaN} in the CBS analytical results because the lecture CBS formulas do not provide a bounded delay for BE traffic.
\end{itemize}

\subsection{Project Structure}

The main components of the implementation are organized into several modules, each responsible for a specific functionality in the simulation framework.
\begin{verbatim}
mini-project-2/
    src/
        main.py
        loader.py
        scheduler.py
        simulation.py
        analysis.py
        model.py
        tools/
            calc_load.py
\end{verbatim}

\begin{itemize}

\item \textbf{model.py}
Defines the core entities used in the system, including \texttt{Link}, \texttt{Node}, \texttt{Stream}, \texttt{Frame}, and \texttt{Route}.

\item \textbf{main.py}
Serves as the entry point of the program and orchestrates the overall execution pipeline of the project.

\item \textbf{loader.py}
Handles the loading and parsing of the test case configuration files.

\item \textbf{simulation.py}
Implements the discrete-event simulator that models frame generation, transmission, propagation, and end-to-end latency measurement.

\item \textbf{scheduler.py}
Implements the scheduling algorithms used in the simulation, including Strict Priority (SP) and Credit-Based Shaper (CBS). The CBS scheduler explicitly tracks credits, blocks AVB queues with negative credit, and resumes transmission when the credit becomes non-negative.

\item \textbf{analysis.py}
Provides analytical methods to compute the Worst-Case Delays (WCDs) for both SP and CBS scheduling. The end-to-end delay is computed as the sum of per-egress-port delays along the route, including the source end-system egress and excluding the final destination node. For the CBS analysis, the reference configuration uses equal-magnitude AVB slopes, i.e. \texttt{idleSlope = 0.5C} and \texttt{sendSlope = -0.5C}.

\item \textbf{calc\_load.py}
Calculates the workload characteristics of each test case, such as the overall AVB traffic load.

\end{itemize}

\subsection{How to Run}
Run the project with:

\begin{verbatim}
cd mini-project-2
python -m src.main <case_id> [duration]
\end{verbatim}

\texttt{case\_id} specifies the test case.
\begin{itemize}
\item 1--3: normal cases
\item 4: starvation of SP case
\end{itemize}

\texttt{duration} is the simulation time in microseconds. The default value is 4000000.0.

For the provided test cases, the implementation assumes a uniform link rate taken from \texttt{topology.default\_bandwidth\_mbps}. In the reference configuration this value is 100 Mbps for every link.

\subsection{Pseudocode}

\textbf{Key Algorithm 1: Worst-Case Delay (WCD) Calculation for CBS}

In the function \texttt{calculate\_wcrt} in \texttt{analysis.py}.

The node delay of each AVB stream is computed using the theoretical formulas from the lecture (\cite{lecture_notes}), considering SPI, HPI, and LPI. The analytical CBS model is only applied to AVB traffic. For BE streams, no analytical CBS bound is reported.

\begin{algorithm}[H]
\caption{Node Delay Calculation for CBS}
\begin{algorithmic}[1]

\State Compute maximum frame sizes for each class
\State $C_A^{max}, C_B^{max}, C_{BE}^{max}$

\State Compute transmission time of current stream
\State $C_i = \frac{size_i \times 8}{linkRate}$

\If{stream class = A}
\State Compute same priority interference
\State $SPI = (\sum C_A - C_i) \times \frac{linkRate}{idleSlope_A}$

\State Compute lower priority interference
\State $LPI = \max(C_B^{max}, C_{BE}^{max})$

\State $nodeDelay = SPI + LPI + C_i$

\ElsIf{stream class = B}
\State Compute same priority interference
\State $SPI = (\sum C_B - C_i) \times \frac{linkRate}{idleSlope_B}$

\State Compute lower priority interference
\State $LPI = C_{BE}^{max}$

\If{$linkRate > idleSlope_A$}
\State $HPI = LPI \times \frac{idleSlope_A}{(linkRate - idleSlope_A)} + C_A^{max}$
\Else
\State System unstable
\EndIf

\State $nodeDelay = SPI + HPI + LPI + C_i$
\EndIf

\end{algorithmic}
\end{algorithm}


\textbf{Key Algorithm 2: WCD Calculation for Strict Priority}

In the function \texttt{calculate\_wcrt\_sp} in \texttt{analysis.py}. 

The analytical worst-case delay under Strict Priority (SP) scheduling is computed using Response Time Analysis (RTA). The delay consists of the frame transmission time, blocking from lower-priority frames, same-priority interference, and interference from higher-priority streams.

\begin{algorithm}[H]
\caption{Worst-Case Delay Calculation for SP}
\begin{algorithmic}[1]

\State $C_i \leftarrow$ transmission time of stream $i$

\State $B_i \leftarrow$ maximum transmission time of lower-priority frames

\State $I_{SP} \leftarrow$ sum of transmission times of same-priority frames

\State Initialize response time:
\State $R \leftarrow C_i + B_i + I_{SP}$

\Repeat
\State $R_{new} \leftarrow C_i + B_i + I_{SP} + \sum_{j \in HP} \left\lceil \frac{R}{T_j} \right\rceil C_j$
\State $R \leftarrow R_{new}$
\Until{$R$ converges}

\State \Return $R$

\end{algorithmic}
\end{algorithm}
